\documentclass[11pt,a4paper]{article}
\usepackage[margin=25mm]{geometry}
\usepackage[T1]{fontenc}
\usepackage{lmodern,microtype,amsmath,amssymb,amsthm,booktabs,graphicx}
\usepackage{xcolor,enumitem,caption}
\usepackage[colorlinks=true,linkcolor=blue!40!black,citecolor=blue!40!black,urlcolor=blue!40!black]{hyperref}
\captionsetup{font=small,labelfont=bf}
\setlist{itemsep=2pt,topsep=4pt,leftmargin=1.6em}
\setlength{\emergencystretch}{2em}
\newtheorem{proposition}{Proposition}
\newcommand{\Sc}{\mathcal{S}}
\title{When a longer route buys time:\\fixed escape plans for regional mine-flood alarms}
\author{}
\date{Working preprint --- October 26, 2026}
\begin{document}
\maketitle
\begin{abstract}
When water floods an underground mine, it spreads through the same tunnels that workers use to escape. Escape routes are planned in advance, and the obvious choice is the shortest way to a shaft. That choice can fail for two reasons. An alarm often tells workers only which part of the mine is flooding, not where the water came in, and it may reach them many minutes after the flood began. By then, a short route that passes near the flooding area may already be under water. We ask whether a longer route, fixed in advance for each starting point and alarm region, can stay open for every possible water source in the region and for every warning delay up to a target, and how much extra walking it costs. For a finite set of flood scenarios, we compute the shortest such route exactly. We apply the method to a 248-km tunnel network taken from a real mine plan, with flooding generated by a simple, uncalibrated model. For a northwest alarm and a 30-minute delay, shortest routes protect 720 of the 1,243 starting points outside the alarm region, and better fixed routes protect 1,200. About half of this gain needs more than 50\% extra walking, and in other regions the gain is smaller or costs more. The method also separates workers who could be saved by a different route from those who would need better information about the source or an earlier warning.
\end{abstract}

\section{Introduction}
A water inrush happens when groundwater, or water stored in old workings, suddenly breaks into an operating mine. It is among the most dangerous accidents underground, because the water and the miners move through the same tunnels. A large mine can have hundreds of kilometers of tunnels connected to the surface by only a few shafts. Water entering at one point runs downhill, fills low sections, and can cut off passages far from where it started. Workers who are still underground have to reach a shaft before the tunnels on their way become too deep to pass.

Mines therefore plan escape routes before any accident happens. Each work area is assigned a route to a shaft, and workers are trained to follow it when the alarm sounds. The natural choice is the shortest route. This paper asks when that choice is wrong, and what a better plan would cost.

The shortest route can fail because, when the alarm sounds, two things are usually unknown. The first is where the water is coming from. A water-level sensor or a report from a work face can tell the control room that water has appeared in some part of the mine, but not necessarily where it broke in. The second is how long the water has already been spreading. Water may travel a long way before it reaches a sensor, and passing the warning on to workers takes more time. A route that is short on the dry map may run close to the flooding area and be closed by the time the worker gets there. A route that is safe for one possible source may fail for another source in the same area.

We ask whether a longer route, chosen in advance, can be safe despite both unknowns. For each starting point and each alarm region, the plan names one route. A worker who hears an alarm for that region sets off on that route immediately. We require the route to reach a shaft for every possible source in the region and for every warning delay up to a target, such as 30 minutes. The target is about how late the worker may set off, not about how long the walk takes. Among the routes that meet this requirement we look for the shortest one, so that the extra walking measures what the protection costs. In this sense a longer route buys time: it lets the plan tolerate a later warning.

To answer the question we need a mine, a way to simulate floods, and an algorithm. Our mine is a tunnel network extracted from a real mine plan, with 1,592 junctions and about 248\,km of tunnels, divided into four alarm regions. For each junction we simulate a flood that starts there, using a simple model of how water volume spreads and when each tunnel becomes impassable. The algorithm is the hard part. Finding the fastest route through one known flood is a standard problem, but here a single route has to work in hundreds of floods at once, and the best route for one source can be the wrong route for another. We solve this problem exactly for every one of the 1,587 starting points, each of the four regions, and three delay targets.

The route choice turns out to matter for large parts of the network, but by very different amounts in different regions. For a northwest alarm and a 30-minute target, the shortest routes protect 720 of the 1,243 starting points outside the region. Choosing routes for this purpose raises the number to 1,200. Only part of that gain is cheap: allowing up to 20\% extra walking protects 812 starting points, and allowing 50\% protects 961. In the southwest, most of the gain needs routes several times longer than the shortest one. The analysis also explains why some starting points cannot be protected by any fixed route. For some of them, at least one possible source leaves no way out at all, so only an earlier warning would help. For others, every source on its own leaves a way out, but the ways out differ and no single route covers them all. Better information about the source could help this second group.

These numbers describe what the planning method can show; they are not safety predictions for this mine. The tunnel geometry is realistic, but the flood model has not been calibrated against measurements, the four regions are simple coordinate quadrants rather than real sensor zones, and walking limits and shaft capacity are not modeled. Section~\ref{sec:limits} discusses these limits. What the study offers is a way to ask, for an escape plan and an alarm system, which workers the plan protects, how much extra walking that protection costs, and whether the remaining failures need a different route, better information, or an earlier warning.

\paragraph{Related work.}
Wu et al.\ \cite{wu2020} find the earliest-arrival escape path through a mine whose tunnels flood over time, assuming the flood is known. A related study places water-level sensors so that an inrush is detected quickly \cite{wu2021}. Our question sits between these two: what should a fixed plan tell workers when detection reports only a region, possibly late? The computational tools come from time-dependent routing and robust optimization. When arriving earlier at a tunnel never means leaving it later, arrival times can be compared directly \cite{kaufman,orda}. Choosing among prepared plans after part of the uncertainty is revealed is studied as $K$-adaptability \cite{hanasusanto}; in our setting, the only thing revealed is the alarm region. The study grew out of a modeling-competition problem on mine water inrush \cite{cumcm}, whose network we use as a second example in Appendix~\ref{app:competition}.

Section~\ref{sec:setup} sets up the model. Section~\ref{sec:compute} explains how the routes are computed. Section~\ref{sec:results} presents the results, and Section~\ref{sec:sensitivity} tests how they change when the inflow or the set of possible sources changes. Section~\ref{sec:limits} discusses what the results do and do not show.

\section{The planning problem}\label{sec:setup}
\subsection{The mine network}
Our test case is a tunnel network extracted from an anonymized CAD plan of a real mine. It has 1,592 junctions and 2,084 tunnels with a total length of about 248\,km. Four junctions at the bottoms of shafts serve as the destinations. Once a worker reaches one of them, we count the escape as complete; the trip up the shaft is not modeled. To build the network we joined nearly touching lines, split lines where they cross, and connected the shafts to the tunnels. The plan gives roof elevations; we treat them as relative floor elevations by assuming all tunnels have the same height. Figure~\ref{fig:geometry} shows the result.

\begin{figure}[tbp]
\centering\includegraphics[width=\linewidth]{figures/cad_geometry_3d.pdf}
\caption{The mine network used in this study. (a) The tunnels in three dimensions, with the vertical scale exaggerated 35 times. (b) Plan view with the four alarm regions and the number of possible sources in each. Stars mark the shafts.}\label{fig:geometry}
\end{figure}

\subsection{What the alarm tells the worker}
We divide the network into four alarm regions, northwest (NW), southwest (SW), southeast (SE) and northeast (NE), at the median $x$ and $y$ coordinates. A real alarm system would define its zones by where its sensors are; quadrants are a simple, reproducible stand-in. Any junction except the four shafts and one auxiliary node can be the place where water enters, which we call the \emph{source}. This gives 1,587 possible sources. Each flood scenario has exactly one source. When the alarm reports region $R$, the worker knows that the source is one of the junctions in $R$, but not which one. Workers can start from any of the same 1,587 junctions.

The plan assigns a route $P(v,R)$ to a worker who starts at junction $v$ and receives an alarm for region $R$. Because the alarm does not give the exact source, the route cannot depend on it. The worker also does not change route on the way. In practice, this is a route map that lists, for each work area, one escape route per alarm region.

\subsection{How late the warning can be}
Let $t_0$ be the time water starts to enter, $t_{\rm det}$ the time the inrush is detected, and $t_{\rm rec}$ the time the worker receives the alarm. A worker who leaves at once therefore sets off after a delay of
\[
 d=t_{\rm rec}-t_0=(t_{\rm det}-t_0)+(t_{\rm rec}-t_{\rm det})
\]
from the start of the inrush. Neither part of this delay is reliably known. Water may travel some distance before it reaches a sensor, so the detection time does not reveal when the inrush started, and passing the warning on takes further time. Any time the worker spends preparing to leave can also be added to $d$.

The planner therefore picks a delay target $\tau$, such as 30 minutes, and asks for a route that works for every delay from $0$ to $\tau$. The worker does not need to know the actual delay. We also assume the alarm reports the correct region; an alarm that can report the wrong region would require a larger set of sources.

\subsection{The flood and walking model}\label{sec:hazard}
To check a route, we need to know, for each source, when each tunnel first gets wet and when it becomes too deep to pass. We compute these times with an event-driven model that tracks the volume of water in every tunnel, inherited from the competition study \cite{cumcm}. Water arriving at a junction divides equally among the level or downhill tunnels leaving it, and goes uphill only when no such tunnel is available. A thin layer of water covers a tunnel before the water moves on beyond it. Dead-end branches fill up, and water then backs up behind them. The depth in a tunnel is its stored volume divided by its floor area. Appendix~\ref{app:flood} gives the full rules.

The main experiments use a constant inflow of $0.5\,\mathrm{m^3/s}$, tunnels 4.99\,m wide and 3\,m high, and a 0.1\,m water layer at the advancing front. A tunnel closes for good when its mean depth reaches 0.3\,m. Workers walk at 1.3\,m/s in dry tunnels. In wet tunnels they walk at 0.7\,m/s in the direction of the water flow and 0.5\,m/s against it. If water reaches a tunnel while a worker is inside it, the worker slows down from that moment. These values define a test case; they are not measured walking speeds or a validated safety threshold.

One property of this model makes the analysis much simpler. Water never recedes, so a worker who enters a tunnel earlier never leaves it later. This is called first-in--first-out (FIFO) travel. It follows that if a route works for a worker who sets off at $\tau$, it also works for any earlier departure. To check a whole range of delays, we only need to check the latest one. Appendix~\ref{app:math} gives the travel-time formula and the proof.

\subsection{What we compute}
Let $\Sc_R$ be the set of possible sources in region $R$, and let $L(P)$ be the length of route $P$. For starting point $v$ and delay target $\tau$, the length of the best plan is
\begin{equation}
\begin{split}
 L_R^*(v,\tau)=\min\bigl\{L(P):{}&P\text{ leads from }v\text{ to a shaft and works}\\
 &\text{for departure at }\tau\text{ for every source in }\Sc_R\bigr\}.
\end{split}
 \label{eq:main}
\end{equation}
If no such route exists, the value is infinite and the starting point cannot be protected. We say a starting point is \emph{protected} when its route works for every source in the region and every delay up to the target. Let $L_0(v)$ be the length of the shortest route from $v$ to a shaft on the dry map. The extra distance that protection requires is
\[
 \rho_R(v,\tau)=\frac{L_R^*(v,\tau)}{L_0(v)}-1,
\]
so a 20\% detour allowance is enough exactly when $\rho_R(v,\tau)\le0.20$. Protection always comes first: we minimize length only among routes that meet the guarantee, so $\rho$ is the least extra walking that the guarantee demands.

We compare this with the \emph{shortest-route plan}, which sends every worker along the shortest dry-map route to the nearest shaft, whatever the alarm region. When several routes tie, a fixed rule picks one. We check the shortest-route plan against exactly the same floods as the rerouted plan. Where two routes are almost equally long, the rerouted plan may simply pick the other one, so a small increase in length does not always mean a real detour.

\section{Finding a route that works for every source}\label{sec:compute}
Computing the best route for each source separately and comparing the answers does not solve our problem: we need one route that works for all sources. So the search has to consider all floods at once. This section explains how we do this exactly for a finite set of sources.

\subsection{One clock for each flood}
Imagine that a partial route carries one clock for each flood scenario, showing when the worker would reach the current junction in that scenario. At each junction the search stores \emph{labels}, each recording a route length and the arrival times in all scenarios. Extending a route along a tunnel updates every clock, using that scenario's wetting time, closure time and walking speed. If the tunnel would already be closed in any scenario when the worker gets there, the extension is discarded.

Two labels at the same junction can be compared. If one is no longer than the other and arrives no later in every scenario, the other label can be dropped: whatever the slower label can do next, the faster one can do too, since arriving earlier never hurts under FIFO travel. If neither is better in every scenario, both are kept. The search expands labels in order of their length plus the remaining dry distance to a shaft, so the first complete route it finds is the shortest route that works in every scenario. A label is also dropped early when a backward calculation shows that, in some scenario, the worker is already too late to reach any shaft.

The number of labels is not capped, so the search is exact. The price is that it can take exponential time on hard instances. Because tunnel lengths are positive and travel is FIFO, an optimal route never needs to visit a junction twice.

\subsection{Adding sources only when needed}
A region has hundreds of possible sources, and carrying hundreds of clocks per label would be slow. Most sources turn out not to matter for a given starting point, so we start with one and add others only when needed:
\begin{enumerate}
\item Start with the source that is worst for the shortest route, and find the shortest route that works for it.
\item Check that route against \emph{every} source in the region, using a fast backward calculation of the latest safe departure time.
\item If some source defeats the route, add that source to the search and solve again. If no source does, accept the route. If no route works even for the sources in the search, report that the starting point cannot be protected.
\end{enumerate}
This procedure is exact; no source is ever sampled away. At each step the search solves an easier version of the full problem, with fewer sources. If the route it finds also passes the check against all sources, no shorter route can work for all of them. If the easier version has no solution, neither does the full problem. Each failed check adds a new source, so the procedure ends after at most as many rounds as there are sources. Appendix~\ref{app:math} states this argument formally.

\subsection{Two ways a fixed route can fail}\label{sec:twofailures}
When no route protects a starting point, there are two possible reasons. In the first, some source leaves the worker no way out at all: even a worker who knew the exact source and could choose the best route for it would not reach a shaft after setting off at $\tau$. In the second, every source on its own leaves a way out, but the ways out go in different directions and no single route works for all of them. We tell the two cases apart with a backward calculation for each single source (Appendix~\ref{app:math}). The distinction matters for what to do next. Better information about the source can help in the second case, but not in the first; there, only an earlier warning would help.

\subsection{Scale of the computation}
We solved the problem for every combination of 1,587 starting points, four regions and three delay targets (10, 20 and 30 minutes), giving 19,044 queries. The two sensitivity studies in Section~\ref{sec:sensitivity} add 19,044 and 6,348 queries. Every one of the 44,436 queries finished within its time limits (three seconds per search and 30 seconds per query), so none had to be reported as unresolved. We fixed the source sets, delay targets, walking speeds and main comparisons before the full results on this network were available; the higher-inflow redesign was added afterwards. We tested the implementation against brute-force enumeration on small networks and checked a sample of routes independently on the mine network. Appendix~\ref{app:checks} reports these tests.

\section{How many workers can a longer route protect?}\label{sec:results}
\subsection{Rerouting protects large parts of the network}
Table~\ref{tab:main} counts, for each alarm region, the starting points outside the region that are protected at a 30-minute target. Workers inside the alarm region, who may be standing next to the source, are discussed in Section~\ref{sec:failures}.

\begin{table}[htbp]
\centering\small
\caption{Starting points outside the alarm region that are protected at a 30-minute delay target. The columns compare the shortest-route plan with rerouted plans that allow up to 20\% or 50\% extra distance, or any length.}\label{tab:main}
\begin{tabular}{@{}lrrrrrrr@{}}\toprule
& & & \multicolumn{4}{c}{Starting points protected} &\\
\cmidrule(lr){4-7}
Region & Sources & Starting points & Shortest route & $\le$+20\% & $\le$+50\% & Any length & Gain\\\midrule
NW & 344 & 1,243 & 720 & 812 & 961 & 1,200 & 480 \\
SW & 448 & 1,139 & 659 & 754 & 760 & 940 & 281 \\
SE & 346 & 1,241 & 1,196 & 1,224 & 1,231 & 1,240 & 44 \\
NE & 449 & 1,138 & 1,075 & 1,079 & 1,092 & 1,110 & 35 \\
\bottomrule

\end{tabular}
\end{table}

The northwest region shows the largest effect. Shortest routes protect 720 of the 1,243 starting points outside it (57.9\%). Rerouting raises this to 1,200 (96.5\%): of the 523 starting points whose shortest route fails, 480 have a longer route that works. The southwest also gains a lot, from 659 to 940 of 1,139 starting points. In the southeast and northeast, shortest routes already protect most starting points, so there is less to gain.

The picture is similar at the shorter delay targets. Rerouting adds 467 protected starting points in the northwest at 10 minutes and 479 at 20 minutes, and 337 and 334 in the southwest (Figure~\ref{fig:detours}b). The gain does not have to shrink as the target grows, because a longer delay lowers both the shortest-route count and the rerouted count.

\begin{figure}[htbp]
\centering\includegraphics[width=\linewidth]{figures/cad_detours.pdf}
\caption{How much rerouting helps starting points outside the alarm region. (a) At a 30-minute target, protection grows with the detour allowance. In the northwest, about half of the gain needs more than 50\% extra distance; in the southwest, most of it needs more than 100\%. (b) With no detour limit, the gain over the shortest-route plan is similar at all three delay targets.}\label{fig:detours}
\end{figure}

The newly protected starting points are not scattered across the network; they form a few connected areas (Figure~\ref{fig:spatial}). The 480 in the northwest form five connected groups, of 402, 59, 17, 1 and 1 junctions. The southwest gains form nine groups, the largest with 203 junctions. This is what we would expect if a few corridors are at risk and rerouting moves everyone who would have used them. It also means the 480 starting points should not be read as 480 independent cases. The CAD plan splits some tunnels into many more junctions than others, so we also weighted each starting point by half the length of the tunnels that meet it. With this weighting, northwest protection rises from 60.0\% to 98.8\% and southwest protection from 52.4\% to 86.0\%, close to the unweighted figures.

\begin{figure}[tbp]
\centering\includegraphics[width=.82\linewidth]{figures/cad_spatial.pdf}
\caption{Starting points outside each alarm region at a 30-minute target. Orange points are protected only by rerouting; red points cannot be protected by any fixed route. The orange points form a few connected areas along shared corridors.}\label{fig:spatial}
\end{figure}

\subsection{What the protection costs}
Much of this gain requires substantial extra walking. In the northwest, a 20\% detour allowance protects 812 starting points, 50\% protects 961, 100\% protects 1,162, and no limit protects 1,200 (Figure~\ref{fig:detours}a). The median newly protected starting point needs 49.0\% extra distance. A study that allowed only small detours would miss about half of the available gain.

The southwest is much more expensive. There, the median newly protected starting point needs 267.4\% extra distance, so its route is about 3.67 times as long as the shortest one. Going from a 100\% allowance to no limit adds 175 starting points. The longest newly protected route from outside an alarm region is about 12.48\,km. Numbers like these show how far rerouting could go in principle. Before any of these routes went into an emergency plan, fatigue, injuries, congestion and shaft capacity would have to be considered, none of which is in our model. A practical plan would add a limit on walking distance or time, and the same method would then show how much protection remains under that limit.

\subsection{An example: 254 extra meters and a different shaft}\label{sec:example}
For a southeast alarm, consider a worker who starts at junction P0036 (Figure~\ref{fig:route}). The shortest route is 6.195\,km long and leads to the shaft at P1081. For at least one of the 346 possible sources, that route is blocked even if the worker sets off at the moment the inrush begins. A route of 6.449\,km to the shaft at P0986 works for all 346 sources at a 30-minute target. The protection costs 254\,m, or 4.1\% extra.

The two routes are almost the same length but have little in common: 6.176\,km of the shortest route and 6.430\,km of the new one are not shared. The new route uses different tunnels and reaches a different shaft that happens to be about as far away. We chose this example to show how rerouting works; how common such improvements are is shown by Table~\ref{tab:main}, not by this example.

\begin{figure}[htbp]
\centering\includegraphics[width=.82\linewidth]{figures/cad_route_example.pdf}
\caption{Two routes from junction P0036 under a southeast alarm. The shortest route (red) is blocked for at least one possible source. The rerouted plan (teal) is only 4.1\% longer, but it follows different tunnels to a different shaft and works for all 346 sources.}\label{fig:route}
\end{figure}

\subsection{Where no fixed route works}\label{sec:failures}
Some starting points cannot be protected even with unlimited detours. For southwest alarms at 30 minutes, there are 199 such starting points outside the region. For 68 of them, at least one source leaves no way out at all, the first kind of failure in Section~\ref{sec:twofailures}. For the other 131, every source on its own leaves a way out, but no single route works for all of them. An alarm system that narrowed down the source could help this second group; for the first group, only an earlier warning would help. We have not designed finer alarm regions or estimated how much a real sensor system would recover.

Workers inside the alarm region fare much worse, since the water may start right next to them. Rerouting raises protection from 14 to 97 of 344 starting points in the northwest, from 31 to 139 of 448 in the southwest, from 25 to 138 of 346 in the southeast, and from 28 to 179 of 449 in the northeast. Appendix~\ref{app:inside} breaks down the remaining failures. Because of this group, the outside-region figures above should not be read as coverage of the whole mine.

\section{How robust are the routes?}\label{sec:sensitivity}
The results so far rest on two assumptions that a planner cannot know precisely: how fast the water comes in, and where it could come from. We change each in turn.

\subsection{Doubling the inflow}
The inflow of $0.5\,\mathrm{m^3/s}$ is an assumption, so we doubled it to $1.0\,\mathrm{m^3/s}$ and asked two questions. Do the routes designed for the original inflow still work? And if they do not, can new routes be found?

\begin{table}[htbp]
\centering\small
\caption{Starting points outside the alarm region protected at a 30-minute target when the inflow is doubled. ``Old routes'' keeps the routes designed for $0.5\,\mathrm{m^3/s}$; ``Rerouted'' designs new routes for $1.0\,\mathrm{m^3/s}$.}\label{tab:flow}
\begin{tabular}{@{}lrrrr@{}}\toprule
& At $0.5\,\mathrm{m^3/s}$ & \multicolumn{3}{c}{At $1.0\,\mathrm{m^3/s}$}\\
\cmidrule(lr){3-5}
Region & Rerouted & Old routes & Shortest route & Rerouted\\\midrule
NW & 1,200 & 855 & 697 & 1,181 \\
SW & 940 & 837 & 632 & 838 \\
SE & 1,240 & 1,235 & 1,194 & 1,238 \\
NE & 1,110 & 1,062 & 1,060 & 1,101 \\
\bottomrule

\end{tabular}
\end{table}

The old routes do not hold up well, but redesigned routes still help a lot (Table~\ref{tab:flow}). In the northwest, only 855 of the 1,200 protected starting points stay protected if the old routes are kept. Redesigning for the higher inflow protects 1,181, against 697 for shortest routes. The northeast shows the same pattern more sharply: only 2 of its 35 newly protected starting points stay protected on their old routes, yet redesign still adds 41 over the shortest routes at the higher inflow. So a plan designed for one inflow may need to be redesigned for another, but rerouting remains worthwhile.

In this model there is a simple way to design for an uncertain inflow. Multiplying a constant inflow by some factor divides every flood event time by the same factor. We used this scaling to produce the doubled-inflow floods, and a direct rerun for the source at junction P0176 reproduced the scaled wetting and closure times exactly. Together with permanent closure and slower walking in wet tunnels, the scaling implies that a route that works at a higher inflow also works at any lower positive inflow, for the same source. A plan designed for an upper bound on the inflow is therefore safe for all smaller inflows, and workers do not need to estimate the inflow while escaping. The argument relies on this scaling property, so it may not hold when water is pumped out, recedes, or follows other hydraulic behavior.

\subsection{A narrower set of possible sources}
Our main source set is deliberately broad: any junction in the reported region could be the source. In reality some locations are much more likely than others. As a contrasting case, we kept only junctions within 50\,m of the water areas marked on the CAD plan, 211 junctions in all. These markings are not a verified risk map; the comparison shows how much the results depend on which sources the plan has to withstand.

\begin{table}[htbp]
\centering\small
\caption{The same comparison as Table~\ref{tab:main} (outside-region starting points, 30-minute target), keeping only sources near the water areas marked on the CAD plan.}\label{tab:near}
\begin{tabular}{@{}lrrrr@{}}\toprule
Region & Sources kept & Shortest route & Rerouted & Gain\\\midrule
NW & 47 & 1,204 & 1,239 & 35 \\
SW & 6 & 1,132 & 1,137 & 5 \\
SE & 103 & 1,241 & 1,241 & 0 \\
NE & 55 & 1,097 & 1,137 & 40 \\
\bottomrule

\end{tabular}
\end{table}

With fewer possible sources, shortest routes already protect almost everyone, and most of the gain disappears (Table~\ref{tab:near}). The northwest gain falls from 480 starting points to 35, the southwest gain from 281 to 5, and the southeast gain to zero. The northeast keeps a gain of 40, but many of these rerouted plans are almost exactly as long as the shortest route; they pick an equally short alternative.

This result does not mean the broad-source results are wrong. Narrowing the source set can only make the best plan safer, but it makes the shortest-route plan safer too, so the advantage of rerouting can shrink. Rerouting is worth most when the alarm leaves the most uncertainty about where the water is. Any claim about its benefit should therefore state which sources the plan is meant to withstand. Choosing a defensible source set for a real mine would need geological and sensor evidence that we do not have.

\section{Limitations and implications}\label{sec:limits}
The results come from a single mine network, so they show that rerouting can matter a great deal, not how much it matters in mines generally. The alarm regions are coordinate quadrants, not real sensor zones. We considered every junction as a source but not points partway along a tunnel, and the shafts themselves were excluded as sources. The counts in our tables are numbers of junctions, which are not the same as numbers of workers or probabilities of survival, and neighboring junctions share corridors.

The flood model is the main physical limitation. It conserves water volume, but it splits flow equally at junctions, spreads water along a tree of tunnels, and tracks only the mean depth of each tunnel. It ignores momentum, changes in depth along a tunnel, losses at junctions, pumping, debris, and the force of moving water on a person. The way it orders simultaneous events can also affect the results. The extracted geometry, the estimated elevations, and the connections to the shafts would all need review by mining engineers. In short, the network makes the routing problem realistically complicated, but it does not make the flood timing realistic.

The model of the workers is equally simple. Routes end at the bottom of a shaft, and the model ignores queues, hoist capacity, fatigue, injury, and the chance of taking a wrong turn. A 12-km route may be mathematically safe and still be a poor emergency instruction. Our unlimited-detour results show how much room for improvement exists before such practical limits are added. A natural next step is to add defensible limits on walking and access and see how much of that room remains.

Within these limits, the analysis suggests a practical way to review an escape plan. First, find the starting points whose shortest routes fail for the information the alarm actually provides. Second, check whether a different fixed route protects them, and how much extra walking it needs. Third, for the points that still fail, ask whether the failure is due to missing information about the source or to a warning that simply comes too late. This turns a single coverage figure into a list of what would have to change.

\section{Conclusion}
A shortest escape route can be the wrong plan when the alarm reports only a region and may arrive late. On a tunnel network taken from a real mine plan, choosing a longer fixed route protected large connected parts of the network that shortest routes left exposed. The extra walking ranged from a few percent to several times the shortest distance, and the gain shrank sharply when fewer sources were considered possible. Two questions remain open: how much of the benefit survives calibrated flood simulations and realistic walking limits, and how alarm regions should be drawn so that fixed routes need smaller detours. Answering them would turn this conditional planning result into evidence for a real emergency plan.

\paragraph{Data, code, and authorship.}
The project archive contains the anonymized network, the source definitions, every individual route result, the solver tests, and the scripts that produce the tables and figures from the stored results (Appendix~\ref{app:flood}). There is no public repository or field dataset. The research, implementation and writing received AI assistance. This follow-up study is separate from the original competition submission and should not be attributed to it.

\begin{thebibliography}{9}\small
\bibitem{wu2020} Q. Wu et al. Finding the earliest arrival path through a time-varying network for evacuation planning of mine water inrush. \emph{Safety Science}, 130:104836, 2020. \href{https://doi.org/10.1016/j.ssci.2020.104836}{doi:10.1016/j.ssci.2020.104836}.
\bibitem{wu2021} Q. Wu, Z. Du, Y. Zhao, H. Xu, and X. Zhang. Optimal location of water level sensors for monitoring mine water inrush based on the set covering model. \emph{Scientific Reports}, 11:2621, 2021. \href{https://doi.org/10.1038/s41598-021-82121-7}{doi:10.1038/s41598-021-82121-7}.
\bibitem{kaufman} D. E. Kaufman and R. L. Smith. Fastest paths in time-dependent networks for intelligent vehicle-highway systems application. \emph{IVHS Journal}, 1(1):1--11, 1993.
\bibitem{orda} A. Orda and R. Rom. Shortest-path and minimum-delay algorithms in networks with time-dependent edge-length. \emph{Journal of the ACM}, 37(3):607--625, 1990.
\bibitem{hanasusanto} G. A. Hanasusanto, D. Kuhn, and W. Wiesemann. $K$-Adaptability in Two-Stage Robust Binary Programming. \emph{Operations Research}, 63(4):877--891, 2015. \href{https://doi.org/10.1287/opre.2015.1392}{doi:10.1287/opre.2015.1392}.
\bibitem{cumcm} China Society for Industrial and Applied Mathematics. \emph{CUMCM 2025, Problem D: Mine water-inrush propagation and escape planning}. Supplied problem statement and network attachments, 2025.
\end{thebibliography}
\clearpage
\appendix
\section{Travel times, guarantees, and exact search}\label{app:math}
Represent each tunnel by two directed arcs. For arc $a$ in source scenario $s$, let $\ell_a>0$ be its length, $w_a^s$ the time it gets wet, $b_a^s$ the time it closes, and $v_a^s$ the walking speed once it is wet. The dry speed satisfies $v_d\ge v_a^s>0$. Ignoring closure, a worker who enters the arc at time $t$ arrives at its far end at
\[
 \widetilde A_a^s(t)=
 \begin{cases}
 t+\ell_a/v_d,&\ell_a\le v_d(w_a^s-t)_+,\\
 \max(t,w_a^s)+[\ell_a-v_d(w_a^s-t)_+]/v_a^s,&\text{otherwise},
 \end{cases}
\]
where $(x)_+=\max(x,0)$. Set $A_a^s(t)=\widetilde A_a^s(t)$ if this is at most $b_a^s$, and $A_a^s(t)=+\infty$ otherwise. Arriving exactly at the closure time counts as success; a stricter safety convention would add an explicit margin.

For a path $P=(a_1,\ldots,a_m)$, composing the arc functions gives $A_P^s(t)$. These functions are nondecreasing (counting $+\infty$ as infeasible), so traversing a path earlier cannot turn a feasible journey into an infeasible one. The latest safe departure along $P$ is
\[
 D_s(P)=\sup\{d\ge0:A_P^s(d)<\infty\},\qquad
 D_R(P)=\min_{s\in\Sc_R}D_s(P),
\]
with $D_s(P)=-\infty$ if no nonnegative departure works. Backward calculations can produce negative deadlines; these only signal that the path is already infeasible at time zero.

\begin{proposition}[Checking one departure time is enough]
For the permanent-closure travel times above, a path $P$ works in every scenario $s\in\Sc_R$ for every departure $d\in[0,\tau]$ if and only if $A_P^s(\tau)<\infty$ for every $s\in\Sc_R$.
\end{proposition}
\begin{proof}
Necessity follows by taking $d=\tau$. For sufficiency, monotonicity gives $A_P^s(d)\le A_P^s(\tau)<\infty$ for every earlier $d$ and every scenario.
\end{proof}

The title's question, how much delay a detour buys, has a matching formulation. For a detour allowance $\varepsilon\ge0$, the longest delay that can be tolerated is
\[
 D_R^*(v,\varepsilon)=\max_{P:L(P)\le(1+\varepsilon)L_0(v)}D_R(P).
\]
With the endpoint convention above, $D_R^*(v,\varepsilon)\ge\tau$ exactly when $L_R^*(v,\tau)\le(1+\varepsilon)L_0(v)$. Fixing the delay target and minimizing length, as in the main text, therefore answers the same question from the other side.

\paragraph{Backward deadlines and the two kinds of failure.}
For a tunnel with parameters $(\ell,w,b,v_d,v_w)$ and a deadline $q$ at its far end, set $e=\min(q,b)$. The latest entry time that reaches the far end by $e$ is
\[
 B(q)=\begin{cases}
 e-\ell/v_w,&\ell\le v_w(e-w)_+,\\
 \min(e,w)-[\ell-v_w(e-w)_+]/v_d,&\text{otherwise}.
 \end{cases}
\]
Applying this map from the last arc back to the first evaluates a fixed path. A reverse label-setting search maximizes it over all routes for one known source; call the result $U_s(v)$. If $\min_s U_s(v)<\tau$, at least one source leaves no way out. If every $U_s(v)\ge\tau$ but no common route exists, the failure comes from incompatible route choices.

\begin{proposition}[Adding sources one at a time]
Fix a starting point $v$, a target $\tau$, and a finite source set $\Sc_R$. Suppose each search either returns an exact shortest route that works for its active subset $T\subseteq\Sc_R$, or certifies that none exists. If a returned route works for all of $\Sc_R$, its length is $L_R^*(v,\tau)$. If the search for some subset finds no route, no route works for $\Sc_R$. Adding a violated source after every failed check reaches one of these conclusions after finitely many additions.
\end{proposition}
\begin{proof}
Every route that works for $\Sc_R$ also works for $T$, so the subset optimum is a lower bound on the full optimum. If the subset route passes all sources, it is also an upper bound, and the two agree. Infeasibility for $T$ implies infeasibility for any larger set. A violated source cannot already be in $T$, so each failed check strictly enlarges $T$, which stays inside the finite set $\Sc_R$.
\end{proof}
The result assumes that each search finishes. A search stopped by a time limit proves nothing about infeasibility; none of our production queries was stopped this way. ``Exact'' refers to the combinatorial search. Floating-point feasibility checks use tolerances of $10^{-7}$ seconds or less, which are not a physical safety margin.

\paragraph{Why choosing a route by the known delay does not help.}
Suppose the worker could pick from several prepared routes according to the exact delay, while the source set stays the same. To cover all delays up to $\tau$, the menu must contain, for delay $\tau$, a route that works for every source. By the first proposition, that route alone also works for every earlier delay. A delay-indexed menu therefore cannot extend the guarantee, although it can save walking when an early alarm allows a shorter route. This argument fails if the observation also narrows the source set, or if the hazard breaks monotone feasibility.

\section{Flood rules and reproducibility}\label{app:flood}
Each tunnel stores a volume between zero and its capacity $WH\ell$, where $W=4.99$\,m and $H=3$\,m. When water first reaches a level or downhill tunnel, the tunnel first receives a layer of volume $Wh_0\ell$ with $h_0=0.1$\,m. If its far junction is still dry, the tunnel then passes water on and joins the propagation tree; if the far junction is already wet, it stores further water. An uphill tunnel fills from below before passing water to a dry junction above it. At a wet junction, incoming flow is divided equally among the level or downhill tunnels still available (with a 0.005\,m elevation tolerance), and uphill tunnels are used only when none remain. When a dead end fills up, the tunnel feeding it stops passing water on and stores it instead. Inflow trapped behind full tunnels is moved through them to the nearest junction that can still store water, and overflow after saturation is counted in the volume balance.

Simultaneous events are ordered by junction and tunnel identifiers. Flow follows the propagation tree, and stored volume never decreases. A tunnel closes when its mean depth reaches 0.3\,m, well below its 3\,m capacity. Walking switches to wet speed when water first enters a tunnel; the model does not track a moving front inside a tunnel. These rules explain both why the model conserves volume and why its schedules should not be read as hydraulic forecasts.

The experiment directory records the source list and input hashes in \texttt{design.json}, and the parameters and analysis design in \texttt{protocol.json}. The files \texttt{full\_*.jsonl} hold all 19,044 primary results, \texttt{near\_*.jsonl} the restricted-source comparison, \texttt{high\_flow/} the doubled-inflow redesign, and \texttt{stress.json} the test of old routes at the higher inflow. The tables and figures are generated from these records by \texttt{make\_cad\_figures.py}; the accompanying README gives the run order and explains which simulation caches must be regenerated. No accident data are used for validation.

\section{Checks on the implementation}\label{app:checks}
We checked the code at three levels. On 100 small random networks with three flood scenarios each, the search agreed with brute-force enumeration of all simple paths in 400 comparisons of feasibility and minimum length, and in a further 100 cases the computed range for the maximum tolerable delay contained the enumerated optimum. On subsets of the mine network, 288 deadline calculations agreed between the single-scenario and many-scenario code, and 72 solutions from the source-adding procedure agreed with a search over all sources at once. Finally, we took five newly protected routes per region, chosen by a fixed ordering, and traced each one forward against every source in its region with separate code; all 7,935 checks passed. Every accepted route that differs from the shortest route was also checked backward against all sources during the optimization.

Across the 1,587 flood simulations, the largest relative error in the volume balance was below $1.02\times10^{-13}$. This confirms the bookkeeping, not the realism of the flood model. The stored results keep, for every query, the source identities, route sequences, the source that decided the outcome, and the solver status.

\section{Starting points inside the alarm region}\label{app:inside}
\begin{table}[htbp]
\centering\small
\caption{Starting points inside the alarm region at a 30-minute target with no detour limit. The last two columns split the unprotected points by the two kinds of failure in Section~\ref{sec:twofailures}: ``No way out'' means at least one source leaves no route at all; ``No common route'' means every source leaves a route but no single route works for all.}
\begin{tabular}{@{}lrrrrr@{}}\toprule
Region & Starting points & Shortest route & Rerouted & No way out & No common route\\\midrule
NW & 344 & 14 & 97 & 172 & 75 \\
SW & 448 & 31 & 139 & 202 & 107 \\
SE & 346 & 25 & 138 & 105 & 103 \\
NE & 449 & 28 & 179 & 215 & 55 \\
\bottomrule

\end{tabular}
\end{table}
For this group, the worker may start at the source itself.

\section{The competition network}\label{app:competition}
This study began with the network from the modeling competition \cite{cumcm}. A first screen used two networks, four regions per network, and eight design sources plus eight held-out sources per region, and picked out 123 starting-point--region pairs for closer study. One promising case on the competition network was then expanded to all 113 possible sources in its southeast quadrant. This was an exploratory investigation, not a confirmation planned in advance.

With the same walking speeds as in this paper, all 548 starting points outside that region were evaluated at three delay targets. At 30 minutes, shortest routes protected 268 of them, a 20\% detour allowance protected 501, and unlimited detours protected 545. So a large rerouting benefit appears in a second network too, although this covers only one region and cannot replace the four-region analysis of the mine network.

The earlier study also showed how a small set of sources can mislead. A route designed for eight sources, which also worked for the eight held-out sources, tolerated a delay of only about 90 seconds once a previously omitted junction was added as a source. Redesigning against the full set fixed that case. We also checked 24 sources at tunnel midpoints for one selected competition route, but this small check does not cover every point along every tunnel of the mine network. These experiences are why the main study enumerates every junction as a possible source.
\end{document}
